% problem-000017 6 四氧化二氮二氧化氮平衡
\section*{6. 四氧化二氮二氧化氮平衡}
\textit{下列为分问。}

\subsection*{6-1}
计算295K下体系达平衡时 $\mathrm { N _ { 2 } O _ { 4 } }$ 和 $\mathrm { N O _ { 2 } }$ 的分压。

\subsection*{6-2}
将上述体系温度升⾄315K，计算达平衡时 $\mathrm { N _ { 2 } O _ { 4 } }$ 和 $\mathrm { N O _ { 2 } }$ 的分压。

\subsection*{6-3}
计算恒压下体系分别在315K和295K达平衡时的体积⽐及物质的量之⽐。

\subsection*{6-4}
保持恒压条件下，不断升⾼温度，体系中NO_{2}分压最⼤值的理论趋近值是多少（不考虑其他反应）？根据平衡关系式给出证明。

\subsection*{6-5}
上述体系在保持恒外压的条件下，温度从295K升⾄315K，下列说法正确的是：

(a) 平衡向左移动 (b) 平衡不移动 (c) 平衡向右移动 (d) 三者均有可能

\subsection*{6-6}
与体系在恒容条件下温度从295K升⾄315K的变化相⽐，恒压下体系温度升⾼，下列说法正确的是(简述理由，不要求计算)：

(a) 平衡移动程度更⼤ (b) 平衡移动程度更⼩ (c) 平衡移动程度不变 (d) 三者均有可能第7题（8分）⼄醇在醋酸菌作⽤下被空⽓氧化是制造醋酸的有效⽅法，然⽽这⼀传统过程远远不能满⾜⼯业的需求。⽬前 ⼯业上多采⽤甲醇和⼀氧化碳反应制备醋酸： $\mathrm { C H _ { 3 } O H + C O  C H _ { 3 } C O O H _ { c } }$ 。第9族元素 $( \mathrm { C o }$

Rh, Ir)的⼀些配合物是上述反应良好的催化剂。以[Rh(CO)_{2}I_{2}]−为催化剂、以碘甲烷为助催化剂合成⼄酸（Monsanto法）的⽰意图如下：

（图：）

\subsection*{7-1}
在催化循环中，A和碘甲烷发⽣氧化加成反应，变为 $\mathbf { B } _ { \circ }$ 。画出B及其⼏何异构体B1的结构⽰意图。

\subsection*{7-2}
分别写出化合物A和D中铑的氧化态及其周围的电⼦数。

\subsection*{7-3}
写出由E⽣成醋酸的反应式（E须⽤结构简式表⽰）。

\subsection*{7-4}
当将上述醋酸合成过程的催化剂改为 $[ \mathrm { I r } ( \mathrm { C O } ) _ { 2 } \mathrm { I } _ { 2 } ] ^ { \cdot }$ −，被称作Cativa法。Cativa法催化循环过程与Monsanto法类似，但中间体C和D（中⼼离⼦均为Ir）有差别，原因在于：由B（中⼼离⼦为Ir）变为 $\mathbf { C } ,$ 发⽣的是CO取代I−的反应；由C到D过程中则发⽣甲基迁移。画出C的⾯式结构⽰意图。
